\documentclass{article}

% ============================================================
% MA 213 In-Class Worksheet Template
% Student-facing worksheet for lecture activities.
% ============================================================

\usepackage[margin=1in]{geometry}
\usepackage{amsmath}
\usepackage{xcolor}
\usepackage{parskip}
\usepackage{array}
\usepackage{enumitem}
\usepackage{booktabs}

\newcommand{\coursenumber}{MA 213}
\newcommand{\weekplaceholder}{10}
\newcommand{\lectureplaceholder}{25}
\newcommand{\lecturetitle}{Setting Up a Hypothesis Test for Paired Differences: Think/Pair/Share}

\newcommand{\nameblank}{\rule{1.75in}{0.4pt}}
\newcommand{\idblank}{\rule{1.55in}{0.4pt}}
\newcommand{\shortblank}{\rule{1in}{0.4pt}}
\newcommand{\longblank}{\rule{0.93\linewidth}{0.4pt}}

\begin{document}

\begin{center}
  {\LARGE \coursenumber{} Week \weekplaceholder{}: Lecture \lectureplaceholder{}}\\
  {\large \lecturetitle{}}
\end{center}

\begin{enumerate}[label=\arabic*.,leftmargin=*,itemsep=.7em]
  \item \textbf{Your name:} \nameblank \hfill \textbf{BU ID:} \idblank
\end{enumerate}

\textbf{Big idea:} 200 students from the High School and Beyond survey each took a reading and a writing test, so the two scores for a given student are \emph{not} independent --- they're paired. We've reduced the two columns to one column of differences (reading $-$ writing) per student. Before the lecture shows you the workflow, work out for yourself how you'd set up the test.

\section*{The Setup}

We have computed \[ \text{diff} = \text{read} - \text{write} \] for each of the 200 students, and we would like to test whether there is a difference between the average reading and writing scores.

\section*{Step 1: Think (2 mins)}

\textbf{(a)} What is the parameter of interest? What is the point estimate?
\vspace{0.7in}

\textbf{(b)} What are the null and alternative hypotheses for testing whether there is a difference between the average reading and writing scores?
\vspace{0.7in}

\textbf{(c)} How does this setting compare to what we've done before (e.g.\ one/two-sample tests for proportions, one-sample test for means)? In particular: how many samples of data do we actually have here?
\vspace{0.7in}

\newpage
\section*{Step 2: Pair (2 mins)}

\begin{enumerate}[label=\arabic*.,leftmargin=*,itemsep=.7em, start=2]
  \item \textbf{Partner's name:} \nameblank \hfill \textbf{BU ID:} \idblank
\end{enumerate}

Compare your answers with your partner. Then discuss:

\begin{itemize}
  \item Did you agree on the parameter of interest and the hypotheses? If not, what was the disagreement?
  \item If the reading and writing scores had been collected from two \emph{different, unrelated} groups of 200 students instead, would the parameter, hypotheses, or analysis change? How?
  \item What test statistic and reference distribution do you expect to use once you compute the standard error of the mean difference?
\end{itemize}

\textbf{Our thoughts:}
\vfill

\end{document}
